% GENERATED FILE. Edit canonical lessons, then run npm run generate:books.
% canonical-source-hash: fnv1a64:5c23f0c58fd24483
% canonical-lessons: MW-C05-hear-mharo, MW-W05-virama, MW-W05-o-matra, MW-C05-mharo, MW-C05-hear-naam, MW-W05-na, MW-C05-naam, MW-C05-hear-hai, MW-W05-ai-matra, MW-C05-answer

\chapter{Mharo Naam: Give One Name}
\label{ch:mw-name-answer}
\hlchaptermodality{5}

\begin{chapteropening}
\textbf{By the end of this chapter:} \emph{I can hear, say, read, and independently write \mw{म्हारो} \mw{नाम} \mw{राम} \mw{है}, then substitute a chosen name in speech.}

A source-attested name statement built by ear before four new signs are isolated and the complete line is retrieved from dictation.
\end{chapteropening}

\section[mhāro]{Hear \emph{mhāro} — one small word for \textquotedblleft{}my\textquotedblright{}}

\label{lesson:MW-C05-hear-mharo}



\begin{quote}
Hear \emph{pāṇī}, say \textbf{water}, say the Marwadi word back, then write its first sign \textbf{\mw{प}} once. Set writing aside. The next chapter begins with the ear.
\end{quote}

\subsection*{What to know first}
\begin{quote}
\emph{mhāro} — \textbf{my}, before a masculine singular noun
\end{quote}

Listen twice: \emph{mhā-ro}. Say it once. The breath after \emph{m} belongs to the word; do not insert a full vowel between \emph{m} and \emph{h}. You do not need to spell it today.

\subsection*{Your turn}
Point to yourself when you hear \emph{mhāro}. Point away when you hear \emph{pāṇī}. Change the order and do it once more.

\begin{tcolorbox}[breakable,colback=teal!4,colframe=teal!35!black,title={Before you move on}]
Say \emph{mhāro} once, then stop.

Sources: \href{https://nepal.sil.org/resources/archives/63768}{SIL International, \emph{Marwari–English Dictionary} (2015)} and \href{https://rsad.artandculture.rajasthan.gov.in/content/dam/doitassets/art-and-culture/Rajasthani-Bhasha-Sahitya-Avm-Sanskriti-Academy-Bikaner/pdf/JJ\_All\_Pdf/JJ\_36\_08\_08\_November\_to\_November\_2007.pdf}{Rajasthan Sahitya Akademi, \emph{Jāgtī Jot}, November 2007}.
\end{tcolorbox}
\section[virāma]{\mw{्} — stop the built-in vowel}

\label{lesson:MW-W05-virama}



\begin{quote}
Write known \textbf{\mw{स}} once, then read \textbf{\mw{म}} as \emph{ma}. Hear \emph{mhāro} once. Today one tiny mark lets \textbf{\mw{म}} close onto the next consonant without \emph{a} between them.
\end{quote}

\subsection*{Script}
\begin{quote}
\mw{्}
\end{quote}

This is the \textbf{virāma}. Under \textbf{\mw{म}}, it removes the built-in \emph{a}: \textbf{\mw{म्}} is \emph{m}, not \emph{ma}. Keep the model visible; no joined form is required yet.

\subsection*{Your turn}
Finger-trace \textbf{\mw{्}} once beneath a printed \textbf{\mw{म}}, then trace one pencil model. Say \textquotedblleft{}vowel off\textquotedblright{} as your hand stops.

\begin{tcolorbox}[breakable,colback=teal!4,colframe=teal!35!black,title={Before you move on}]
Cover the mark and draw it once beneath \textbf{\mw{म}}.
\end{tcolorbox}
\section[o-mātrā]{\mw{ो} — turn \emph{ra} into \emph{ro}}

\label{lesson:MW-W05-o-matra}



\begin{quote}
Write known \textbf{\mw{सा}} once and point to \textbf{\mw{ी}} in \textbf{\mw{पाणी}}. Write \textbf{\mw{र}} once, then draw yesterday's vowel-off mark beneath a spare \textbf{\mw{म}}. Today adds one sign to the right of \textbf{\mw{र}}.
\end{quote}

\subsection*{Script}
\begin{quote}
\mw{र} + \mw{ो} $\to$ \mw{रो}
\end{quote}

Attached \textbf{\mw{ो}} replaces the built-in \emph{a}. Read \textbf{\mw{रो}} as \emph{ro}. It is one sign, even though it reaches above and to the right.

\subsection*{Your turn}
Keep \textbf{\mw{रो}} visible and copy it once. Point to \textbf{\mw{र}}, then \textbf{\mw{ो}}, and read \emph{ro}.

\begin{tcolorbox}[breakable,colback=teal!4,colframe=teal!35!black,title={Before you move on}]
Cover the model, wait three seconds, and write \textbf{\mw{रो}} once.
\end{tcolorbox}
\section[mhāro]{\mw{म्हारो} — the heard word becomes readable}

\label{lesson:MW-C05-mharo}



\begin{quote}
Say the Marwadi word for \textbf{my} and recall what respectful \textbf{\mw{सा}} does. From the heard cue \emph{pāṇī}, write \textbf{\mw{पाणी}}. Then write the vowel-off mark and \textbf{\mw{रो}} from memory.
\end{quote}

\subsection*{What to know first}
\begin{quote}
\textbf{\mw{म्}} + \textbf{\mw{हा}} + \textbf{\mw{रो}} $\to$ \textbf{\mw{म्हारो}} — \emph{mhāro} — my
\end{quote}

Every piece is known. The virāma removes the vowel from \textbf{\mw{म}} so \textbf{\mw{म्ह}} begins \emph{mh}, not \emph{maha}.

\subsection*{Your turn}
Look for five seconds, cover \textbf{\mw{म्हारो}}, wait five seconds, then write it. Uncover and compare \textbf{\mw{म्ह}}, \textbf{\mw{ा}}, and \textbf{\mw{रो}} separately. Repair only the piece that differs.

\begin{tcolorbox}[breakable,colback=teal!4,colframe=teal!35!black,title={Before you move on}]
Read your repaired word once without romanization.

Sources: \href{https://nepal.sil.org/resources/archives/63768}{SIL International, \emph{Marwari–English Dictionary} (2015)} and \href{https://rsad.artandculture.rajasthan.gov.in/content/dam/doitassets/art-and-culture/Rajasthani-Bhasha-Sahitya-Avm-Sanskriti-Academy-Bikaner/pdf/JJ\_All\_Pdf/JJ\_36\_08\_08\_November\_to\_November\_2007.pdf}{Rajasthan Sahitya Akademi, \emph{Jāgtī Jot}, November 2007}.
\end{tcolorbox}
\section[nām]{Hear \emph{nām} — name}

\label{lesson:MW-C05-hear-naam}



\begin{quote}
Hear \textbf{\mw{राम}-\mw{राम} \mw{सा}} and answer it. Retrieve \textbf{\mw{पाणी}} once by ear, voice, eye, and writing. Then hear \emph{mhāro}, say \textbf{my}, and read \textbf{\mw{म्हारो}} once.
\end{quote}

\subsection*{What to know first}
\begin{quote}
\emph{nām} — name
\end{quote}

Listen twice, then say it once. Marwadi \emph{nām}, Sanskrit \emph{nāman}, and English \emph{name} are distant members of the same old word family. That history is a memory hook; today's task is only sound and meaning.

\subsection*{Your turn}
Touch your name on a card when you hear \emph{nām}. Touch yourself when you hear \emph{mhāro}. Change the order once.

\begin{tcolorbox}[breakable,colback=teal!4,colframe=teal!35!black,title={Before you move on}]
Say \emph{mhāro nām} and give the two-word meaning: \textbf{my name}.

Sources: \href{https://nepal.sil.org/resources/archives/63768}{SIL International, \emph{Marwari–English Dictionary} (2015)} and \href{https://dsal.uchicago.edu/dictionaries/soas/}{Turner, \emph{A Comparative Dictionary of the Indo-Aryan Languages}}.
\end{tcolorbox}
\section[na]{\mw{न} — a front-of-the-mouth \emph{n}}

\label{lesson:MW-W05-na}



\begin{quote}
Write independent \textbf{\mw{आ}} once. Say \emph{nām} and its meaning. Read known \textbf{\mw{ण}} with the tongue curled back. Today the tongue moves forward.
\end{quote}

\subsection*{Script}
\begin{quote}
\mw{न}
\end{quote}

Read \textbf{\mw{न}} as \emph{na}, with the tongue near the teeth. Do not replace it with retroflex \textbf{\mw{ण}} from \textbf{\mw{पाणी}}.

\subsection*{Your turn}
Keep \textbf{\mw{न}} visible. Finger-trace it once, then trace one pencil model. Point to \textbf{\mw{न}} when you hear \emph{na} and \textbf{\mw{ण}} when you hear retroflex \emph{ṇa}.

\begin{tcolorbox}[breakable,colback=teal!4,colframe=teal!35!black,title={Before you move on}]
Cover the model and write \textbf{\mw{न}} once.
\end{tcolorbox}
\section[nām]{\mw{नाम} — name, with no new sign}

\label{lesson:MW-C05-naam}



\begin{quote}
Write known \textbf{\mw{भ}} and the vowel-off mark \textbf{\mw{्}} once each. Say \emph{nām} and its meaning, then write \textbf{\mw{न}} once.
\end{quote}

\subsection*{What to know first}
\begin{quote}
\textbf{\mw{न}} + \textbf{\mw{ा}} + \textbf{\mw{म}} $\to$ \textbf{\mw{नाम}} — \emph{nām} — name
\end{quote}

Nothing new is hiding here: dental \textbf{\mw{न}}, known long \textbf{\mw{ा}}, and known \textbf{\mw{म}}.

\subsection*{Your turn}
Look for five seconds, cover \textbf{\mw{नाम}}, wait five seconds, and write it once. Uncover, compare one piece at a time, and repair only a differing piece.

\begin{tcolorbox}[breakable,colback=teal!4,colframe=teal!35!black,title={Before you move on}]
Read \textbf{\mw{म्हारो} \mw{नाम}} and say \textbf{my name}.
\end{tcolorbox}
\section[hai]{Hear \emph{hai} — the small closing word}

\label{lesson:MW-C05-hear-hai}



\begin{quote}
Read and write \textbf{\mw{आभार}} once. Hear \emph{mhāro} and \emph{nām} in either order, then give each meaning.
\end{quote}

\subsection*{What to know first}
\begin{quote}
\emph{hai} — is
\end{quote}

Listen twice, then say it once. This tiny copula closes the name statement. Do not spell it yet; one vowel sign is still missing from your hand.

\subsection*{Your turn}
Hear \emph{mhāro nām rām hai}. Tap four beats: \textbf{my | name | Ram | is}. English changes the order in translation: \textbf{My name is Ram.}

\begin{tcolorbox}[breakable,colback=teal!4,colframe=teal!35!black,title={Before you move on}]
Say the four-beat line once from sound: \emph{mhāro nām rām hai}.

Sources: \href{https://www.rgmcet.edu.in/assets/img/documents/Campus\%20Life/Other\%20Reports.pdf}{Ek Bharat Shreshtha Bharat language report} and \href{https://rsad.artandculture.rajasthan.gov.in/content/dam/doitassets/art-and-culture/Rajasthani-Bhasha-Sahitya-Avm-Sanskriti-Academy-Bikaner/pdf/JJ\_All\_Pdf/JJ\_36\_08\_08\_November\_to\_November\_2007.pdf}{Rajasthan Sahitya Akademi, \emph{Jāgtī Jot}, November 2007}.
\end{tcolorbox}
\section[ai-mātrā]{\mw{ै} — turn \emph{ha} into \emph{hai}}

\label{lesson:MW-W05-ai-matra}



\begin{quote}
Say \textbf{\mw{आभार}} for formal thanks. Say the closing word \emph{hai}, then write known \textbf{\mw{ह}} once.
\end{quote}

\subsection*{Script}
\begin{quote}
\mw{ह} + \mw{ै} $\to$ \mw{है}
\end{quote}

Attached \textbf{\mw{ै}} replaces the built-in \emph{a}. Read the whole word \textbf{\mw{है}} as \emph{hai}. Do not confuse \textbf{\mw{ै}} with the taller \textbf{\mw{ो}} in \textbf{\mw{रो}}.

\subsection*{Your turn}
Keep \textbf{\mw{है}} visible and copy it once. Point to \textbf{\mw{ह}}, then \textbf{\mw{ै}}, and read \emph{hai}.

\begin{tcolorbox}[breakable,colback=teal!4,colframe=teal!35!black,title={Before you move on}]
Cover the model, wait three seconds, and write \textbf{\mw{है}} once.
\end{tcolorbox}
\section[mhāro nām rām hai]{\mw{म्हारो} \mw{नाम} \mw{राम} \mw{है।} — say one name}

\label{lesson:MW-C05-answer}



\begin{quote}
Hear \emph{mhāro, nām, rām, hai} in mixed order. Give each meaning.
\end{quote}

\subsection*{What to know first}
\begin{quote}
\textbf{\mw{म्हारो} | \mw{नाम} | \mw{राम} | \mw{है।}}
\end{quote}

>

\begin{quote}
\emph{mhāro nām rām hai} — My name is Ram.
\end{quote}

The written line adds no new sign. \textbf{\mw{है}} is known \textbf{\mw{ह}} plus \textbf{\mw{ै}}. Replace \textbf{\mw{राम}} with your chosen name when speaking; keep Ram for the first writing attempt because every sign is already yours.

\subsection*{Your turn}
1. hear the line and point to the person naming themself 2. say the line once with \emph{rām} 3. read the four groups without romanization

\begin{tcolorbox}[breakable,colback=teal!4,colframe=teal!35!black,title={Before you move on}]
Cover every written model. Hear the line once and write it. Then uncover and repair one group, not the whole sentence.

Sources: \href{https://www.rgmcet.edu.in/assets/img/documents/Campus\%20Life/Other\%20Reports.pdf}{Ek Bharat Shreshtha Bharat language report} and \href{https://rsad.artandculture.rajasthan.gov.in/content/dam/doitassets/art-and-culture/Rajasthani-Bhasha-Sahitya-Avm-Sanskriti-Academy-Bikaner/pdf/JJ\_All\_Pdf/JJ\_36\_08\_08\_November\_to\_November\_2007.pdf}{Rajasthan Sahitya Akademi, \emph{Jāgtī Jot}, November 2007}.
\end{tcolorbox}
